\section{The data, prime sums, heights and twisted lattices}\label{sec:setup}

This section fixes the objects that depend on the length $x$ of the prime sum, defines the
quantities that the main inequality compares, and proves the general lower bound for a Gram
determinant on which Theorem~C rests. The Lean files are those of
\leanfile{New/Reduction/Data}, \leanfile{New/Reduction/PrimeSums} and
\leanfile{New/Reduction/Lattices}.

\subsection{Parameters, the field and its places}

\begin{definition}[\lean{New.Parameters}]\label{setup:def:parameters}
The \emph{parameters} are a positive integer $q_0$ divisible by $18480 = 16\cdot3\cdot5\cdot7\cdot11$
and a prime $d$ with $q_0 \mid 2(d+1)$ and $q_0 \nmid d+1$; equivalently
$d \equiv q_0/2 - 1 \bmod q_0$. Put $K = \Q(\sqrt{-d})$ and let $\chiK$ be its quadratic
character.
\end{definition}

Then $d \equiv 7 \bmod 8$ (\lean{New.Parameters.instPrimeSevenModEight}), every prime $p$ dividing
$q_0$ (or $q_1$ below) has $\chiK(p) = 1$, so splits in $K$
(\lean{New.Parameters.chiK_eq_one_of_dvd}), and a prime that does not split in $K$ is at least
$13$ (\lean{New.Parameters.thirteen_le_of_chiK_ne_one}).

For an integer $s \ge 1$ let $l_0$ be the least power of $3$ above $27s$, and put
\[
  T = 3\, l_0\, 2^s, \qquad q_1 = q_0 T, \qquad q_2 = (q_0 d)^{s+6}
\]
(\lean{New.towerDeg}, \lean{New.Parameters.q₁}, \lean{New.Parameters.q₂}). Let $E = K(\zeta)$ with
$\zeta$ a primitive $q_1$-th root of unity (\lean{New.Parameters.E}). The cyclotomic polynomial of
order $q_1$ is irreducible over $K$ (\lean{New.Parameters.irreducible_cyclotomic}), so
$\Gal(E/K) \cong (\Z/q_1)^\times$; the field $E$ is a CM field, abelian over $\Q$. Each infinite
place of $E$ is complex, and a fixed embedding $E \to \C$ over $K$ composed with the elements of
$\Gal(E/K)$ gives one embedding for each place. So the set $\Phi$ of infinite places of $E$ is
identified with the group $(\Z/q_1)^\times$ (\lean{New.Parameters.placeMulEquiv}), and these
embeddings form a CM type. We write $f = \#\Phi = T\varphi(q_0)$, and for an integer $\nu$ prime to
$q_1$ and $\tau \in \Phi$ we write $\nu\tau$ for the product of $\tau$ with the place of the class
of $\nu$. A place has a \emph{class modulo $q_0$}, its image in $(\Z/q_0)^\times$; exactly $T$
places have class $1$ (\lean{New.Parameters.card_placeClass_eq_one_eq_towerDeg}). Finally
\[
  \Omega = \exp\bigl(\gamma - \Re \tfrac{L'}{L}(1,\chiK)\bigr)\cdot \frac{2\pi}{\rdE},
\]
with $\gamma$ Euler's constant and $\rdE$ the root discriminant of $E$
(\lean{New.Parameters.bigOmega}).

\subsection{Admissible data}

\begin{definition}[\lean{New.Data}]\label{setup:def:data}
\emph{Admissible data} $D$ for the parameters $(q_0, d)$ consist of the following.
\begin{enumerate}[label=(\alph*)]
\item \emph{Scales}: integers $x$, $x_1$, $b$ with
  $(\log b)^{100} \le x$, $x^{10^6} \le x_1$, $x_1^{10^6} \le b$, $(\log x_1)^6 \le \log b$ and
  $\log\log b \ge 10^{12} c_P$, where $c_P$ is a constant that depends on $q_0$ and $d$ only
  (\lean{New.Parameters.constantBound}). Put $a = x b$.
\item \emph{Ranks}: a power of two $m \in [\sqrt{x}, 2\sqrt{x}]$, the number of lifts; a power of
  three $r \in [x^5, 3x^5]$; and the exponent $s$ with $2^s = m^{22}$, which fixes $E$, $q_1$ and
  $q_2$ as above, with $s \le 32 \log x$. Further $q_1 x_1$ divides $b$. Put
  \[
    k = r + 1 - m, \qquad t = 5r, \qquad n_j = k + j(t+1), \qquad \Delta = 2r^8 .
  \]
\item \emph{Places}: the set $\Phi_a$ of \emph{$a$-places}, the $T$ places of class $1$ modulo
  $q_0$, and a set $\Phi_b$ of \emph{$b$-places}, all of class $-1$ modulo $q_0$, with
  $\#\Phi_b \ge T/\log b$, that is \emph{sparse}: for every prime $\ell \le kt$ and all
  $0 < i < j < l_0$ there are no $\sigma, \tau, \rho \in \Phi_b$ with
  $\tau/\sigma \equiv \pm\ell^i$ and $\rho/\sigma \equiv \pm\ell^j$ (independent signs) modulo the
  part of $q_1$ prime to $\ell$ (\lean{New.Parameters.IsSparse}).
\item \emph{Weights}: $\kappa(\tau) = b = a/x$ for $\tau \in \Phi_b$ and $\kappa(\tau) = a/x_1$
  for all other places; and $\kappa_j(\tau) = \kappa(\tau) + (j + \tfrac12)\Delta$, the weight of
  the $j$-th lift.
\item \emph{Primes}: distinct primes $p_0, p_1 \in [b^{10}, b^{11}]$ and a prime $N$, all
  $\equiv 1 \bmod q_2^2$, with an integer of $K$ of norm $M = p_0 p_1$ (its conjugate is written
  $\pi_M$) and an integer $\pi_N$ of $E$ with $\pi_N \overline{\pi_N} = N$. The level $N$ satisfies
  $\log N_- \le \log N \le \log N_- + \sqrt{\log b}$, where
  \[
    N_- = \bar\kappa^{\,m}\, \Bigl( \frac{M\,\rdE}{2\pi} \Bigr)^{m-1} b^{-15/2}
  \]
  and $\bar\kappa$ is the geometric mean of the weights $\kappa(\tau)$ over the $f$ places.
\end{enumerate}
\end{definition}

Table~\ref{setup:tab:sizes} lists the sizes. All of them follow from the definition
(\leanfile{New/Reduction/Data/Sizes}, \leanfile{New/Reduction/Data/Magnitudes},
\leanfile{New/Reduction/Data/LevelAndTwists}; for $\rdE$, \lean{New.Data.rootDisc_mem_Icc}). In the Lean text
each quantity carries an upper and a lower bound for its logarithm, and an inequality between two
quantities is proved by comparing these bounds, which is a linear problem in $\log x$, $\log x_1$,
$\log b$, $\log\log b$, $\log\log x_1$ and $c_P$ (\lean{New.Data.le_of_hasLogBound}). We use such
inequalities freely and cite them by this name.

\begin{table}[ht]
\centering
\begin{tabular}{lll}
\toprule
quantity & role & size \\
\midrule
$\log x_1$, $\log b$ & scales & $(\log x_1)^6 \le \log b$, $(\log b)^{100} \le x$ \\
$m$ & number of lifts & $[\sqrt{x}, 2\sqrt{x}]$ \\
$r$, $k$, $t$ & ranks & $r \in [x^5, 3x^5]$, $r/2 \le k \le r$, $t = 5r$ \\
$s$, $l_0$ & tower & $s \le 32\log x$, $27s < l_0 \le 81s$ \\
$T$, $f$, $q_1$ & places & $x^{11} \le m^{22} = 2^s \le T \le f \le q_1 \le x^{12}$ \\
$\rdE$, $1/\Omega$ & field & $[x^{10}, x^{12}]$, $[1, x^{13}]$ \\
$\Delta$ & gap of weights & $2r^8$ \\
$a$, $\kappa(\tau)$ & orders & $a = xb$, $\ a/x_1 \le \kappa(\tau) \le b$ \\
$M = p_0p_1$ & principal level & $[b^{20}, b^{22}]$ \\
$N$ & level & $\log N \in [m\log b, 30\, m\log b]$ \\
$f a \log^3 b$ & unit of the errors & \lean{New.Data.errorUnit} \\
\bottomrule
\end{tabular}
\caption{Sizes of admissible data. In increasing order: $\log b$, $x$, $x_1$, $b$, $M$, $N$.}
\label{setup:tab:sizes}
\end{table}

\begin{theorem}[Existence of the data; \lean{New.eventually_exists_data}]\label{setup:thm:existence}
Let $0 < \eps < 1/100$. For every large integer $x$ there are admissible data with first scale
$x$, with $x_1 = \floor{\exp(x^{\eps/6})}$ and with $x^\eps \le \log b \le x^\eps + 1$.
\end{theorem}

The proof is a construction (\leanfile{New/Reduction/Data/Existence}). The $b$-places are the
places $-g^e$, $g = 1 + 3l_0q_0$, with $e$ in a Behrend-type set without solutions of several
linear equations at once (\lean{New.exists_three_term_free}; Behrend, 1946; Shapira, 2006): a
forbidden triple of places would give a solution. The primes come from the prime number theorem in progressions, with the logarithm of
the prime in a prescribed window (\leanfile{Lib/Analytic/Primes}); the integer of norm $M$ from a
pigeonhole in the class group of $K$; $\pi_N$ from $N \equiv 1 \bmod 4$, a sum of two squares.

\subsection{Prime sums}

\begin{definition}[\lean{New.Data.primeSum}, \lean{New.Data.jetSize},
\lean{New.Data.directionAmount}, \lean{New.Data.mainAmount}, \lean{New.Data.mainPrimeSum},
\lean{New.Data.totalRemainder}]\label{setup:def:primesums}
For real $y \ge 0$, a place $\tau$ and an integer $j \ge 0$ put
\[
  \sfin(y,\tau) = \sum_{1 \le \nu \le a} \LambdaK(\nu)\,
     \Bigl( \frac{y}{\nu} - \kappa(-\nu\tau) \Bigr)_{\!+},
  \qquad
  \sinf(j) = \log j! + j \log\frac{M}{\Omega},
\]
where $u_+ = \max(u,0)$. (If $\nu$ is not prime to $q_1$, $-\nu\tau$ is read as $\tau$; then
$\LambdaK(\nu) = 0$ anyway, as $\Lambda$ lives on prime powers and $\chiK(p) = 1$ for $p \mid q_1$.) The \emph{direction amount}, the \emph{main amount}, the \emph{expected value} and
the \emph{total remainder} are
\begin{align*}
  \hdir &= \sum_{\tau \in \Phi} \bigl( \sinf(\kappa(\tau)) - \sfin(\kappa(\tau),\tau) \bigr),
  & \hmain &= m\,k\,\hdir, \\
  \smain &= a\,\Bigl( \log a - \frac1f \sum_{\tau\in\Phi} \log\kappa(\tau) + c_1 \Bigr),
  & \rem &= \sum_{\tau \in \Phi_a} \bigl( \sfin(a,\tau) - \smain \bigr),
\end{align*}
with $c_1 = \Re G_1(1)$ as in Definition~\ref{zeros:def:remainder}.
\end{definition}

The quantity $\sinf(j)$ is the logarithmic size charged to a derivative of order $j$ of a theta
function, and $\sfin(y,\tau)$ is the sum of $2v_p\log p$ over the primes $p$ that do not split in $K$,
with $p^{v_p}$ the predicted divisibility of such a derivative (Section~\ref{sec:der}). The main amount is
large and positive: $\hmain \ge x^4 f a$ (\lean{New.Data.le_mainAmount}). It will cancel. What
does not cancel is the total remainder, and it is a combination of the remainders of
Section~\ref{sec:zeros}:

\begin{proposition}[\lean{New.Data.totalRemainder_eq}]\label{setup:prop:rem}
\[
  \rem = a\,\Bigl( \#\Phi_b\,\bigl( R_{q_0}(x) - R_{q_0}(x_1) \bigr) + \#\Phi_a\, R_1(x_1) \Bigr).
\]
\end{proposition}

\begin{proof}
Fix $\nu \le a$ prime to $q_1$ (for other $\nu$ both sides of the next display vanish) and let
$\tau$ run over $\Phi_a$. Then $-\nu\tau$ runs once over
the places of the class of $-\nu$ modulo $q_0$. The $b$-places have class $-1$. So $-\nu\tau$ is a
$b$-place for exactly $\#\Phi_b$ values of $\tau$ if $\nu \equiv 1 \bmod q_0$, and for none
otherwise (\lean{New.Data.card_filter_mem_bPlaces}). As $\kappa = a/x$ on $\Phi_b$ and
$\kappa = a/x_1$ elsewhere,
\[
  \sum_{\tau\in\Phi_a} \LambdaK(\nu)\Bigl(\frac a\nu - \kappa(-\nu\tau)\Bigr)_{\!+}
  = \#\Phi_a\,\lambda_1(\nu)\Bigl(\frac a\nu - \frac a{x_1}\Bigr)_{\!+}
  + \#\Phi_b\,\lambda_{q_0}(\nu)\Bigl[\Bigl(\frac a\nu - \frac a{x}\Bigr)_{\!+}
       - \Bigl(\frac a\nu - \frac a{x_1}\Bigr)_{\!+}\Bigr].
\]
For $0 < z \le a$, the definition \eqref{zeros:eq:R} gives
$\sum_{\nu \le a} \lambda_q(\nu)(a/\nu - a/z)_+ = a\,(R_q(z) + \log z/\varphi(q) + c_q)$
(\lean{New.sum_lambdaKResidueClass_mul_posPart}). Summing over $\nu$,
\[
  \sum_{\tau\in\Phi_a} \sfin(a,\tau)
  = a\,\Bigl( \#\Phi_a\bigl(R_1(x_1) + \log x_1 + c_1\bigr)
     + \#\Phi_b\Bigl(R_{q_0}(x) - R_{q_0}(x_1) + \frac{\log(x/x_1)}{\varphi(q_0)}\Bigr) \Bigr);
\]
the constant $c_{q_0}$ has cancelled. On the other hand
$\sum_\tau \log\kappa(\tau) = f\log(a/x_1) - \#\Phi_b\log(x/x_1)$
(\lean{New.Data.sum_log_weight}) and $f = \#\Phi_a\,\varphi(q_0)$, so $\#\Phi_a\,\smain = a\,(\#\Phi_a(\log x_1 + c_1) + \#\Phi_b\log(x/x_1)/\varphi(q_0))$
(\lean{New.Data.card_aPlaces_mul_mainPrimeSum}). Subtract.
\end{proof}

Hence the two kinds of places: only $\nu \equiv 1 \bmod q_0$ moves an $a$-place to a $b$-place,
where the weight is larger, so the class modulo $q_0$, and with it the Dirichlet characters
modulo $q_0$, enters through the weights alone.

\subsection{Heights}

Let $\mathcal{I}$ be the set of pairs $\iota = (\ell, \iota)$ of a prime $\ell$ and a field
isomorphism $\iota \colon \C \isom \Qlbar$ (\lean{New.PrimePadicIso}); we write $\iota \mid \ell$,
and $\norm{\cdot}_\ell$ for the absolute value of $\Qlbar$ with $\norm{\ell}_\ell = \ell^{-1}$. A
complex number $z$ is \emph{integral at $\iota$} if $\norm{\iota z}_\ell \le 1$. The conjugate
isomorphism is $\bar\iota(z) = \iota(\bar z)$.

\begin{definition}[\lean{New.HeightLe}]\label{setup:def:height}
A family $y = (y_\iota)_{\iota \in \mathcal{I}}$ of complex numbers has \emph{height at most $c$}
if there are a complex number $H \ne 0$, a real number $\alpha$ and real numbers $h_\ell \ge 0$
such that
\[
  \abs{\sigma H} \le e^{\alpha} \ \text{ for all } \sigma \in \Aut(\C), \qquad
  \norm{\iota(H y_\iota)}_\ell \le \ell^{h_\ell} \ \text{ for all } \iota \mid \ell, \qquad
  \alpha + \sum_{\ell \in S} h_\ell \log \ell \le c
\]
for every finite set $S$ of primes.
\end{definition}

The multiplier $H$, and with it the bound at infinity, is the same for all $\iota$, while the
number $y_\iota$ may depend on $\iota$. This is what later lets each prime use its own twist.
Heights add under pointwise products and powers (\lean{New.HeightLe.mul}, \lean{New.HeightLe.pow})
and are monotone (\lean{New.HeightLe.mono}); they may be negative.

\begin{lemma}[Product formula; \lean{New.HeightLe.neg_le_of_norm_map_le}]\label{setup:lem:product}
Let $z \ne 0$ be a complex number. If the constant family $z$ has height at most $c$, and
$\abs{\sigma z} \le e^{c'}$ for all $\sigma \in \Aut(\C)$, then $c + c' \ge 0$.
\end{lemma}

\begin{proof}
Take $H$, $\alpha$, $h_\ell$ as in the definition and put $w = Hz \ne 0$. All conjugates
$\sigma w$ are bounded by $e^{\alpha + c'}$, and $\norm{\iota w}_2 \le 2^{h_2}$ for all $\iota \mid 2$;
the latter makes $w$ algebraic (\lean{New.isAlgebraic_of_norm_complexToCp_le}), of some degree
$n$. Its norm
$\Nm(w) \in \Q^\times$ has $\abs{\Nm(w)} \le e^{n(\alpha + c')}$ and
$\norm{\Nm(w)}_\ell \le \ell^{n h_\ell}$ for each prime $\ell$
(\lean{New.HeightLe.norm_algebraMap_norm_gen_le}). By the product formula for $\Q$,
$0 = \log\abs{\Nm(w)} + \sum_\ell \log\norm{\Nm(w)}_\ell \le n(\alpha + c') + n\sum_\ell h_\ell\log\ell
\le n(c + c')$, the sums being over the finitely many $\ell$ with $\norm{\Nm(w)}_\ell \ne 1$.
\end{proof}

\subsection{Twisted lattices}

\begin{definition}[\lean{New.TwistedLattice}]\label{setup:def:lattice}
Let $V$ be a complex vector space and $\mathcal{T}$ a set with a distinguished element $1$, whose
elements are called \emph{twists}; inner products are linear in the second variable. A \emph{twisted lattice} on $V$ consists of
\begin{itemize}
\item for each twist $\xi$ a complex inner product space $U_\xi$ and a linear map
  $\Theta_\xi \colon V \to U_\xi$;
\item for each prime $\ell$ a set of twists \emph{allowed at $\ell$};
\item for each $\xi$ and each $\iota \mid \ell$ a module $U_\xi(\iota) \subset U_\xi$ over the
  ring $\cO_\iota$ of complex numbers integral at $\iota$, the vectors \emph{integral at $\iota$},
  such that $\pair{\varphi, \varphi'}$ is integral at $\iota$ whenever $\xi$ is allowed at $\ell$,
  $\varphi \in U_\xi(\bar\iota)$ and $\varphi' \in U_\xi(\iota)$;
\item finitely many linear forms on $U_1$, the \emph{coordinates}, in which the inner product has a
  rational matrix that is positive definite on the whole coordinate space.
\end{itemize}
A vector $v \in V$ is \emph{integral at $\iota \mid \ell$} if $\Theta_\xi v \in U_\xi(\iota)$ for
every twist $\xi$ allowed at $\ell$ (\lean{New.TwistedLattice.integralVectors}).
\end{definition}

For $\sigma \in \Aut(\C)$ and $v \in V$ let $\sigma\Theta_1 v$ be the coordinate vector of
$\Theta_1 v$ with $\sigma$ applied to each coordinate, and pair such vectors by the rational
matrix. Write $\norm{\Theta_1 v}_\sigma^2$ for the product of $\sigma\Theta_1 v$ with itself
(\lean{New.TwistedLattice.normSqAt}).

\begin{definition}[\lean{New.TwistedLattice.IndexLe}, \lean{New.TwistedLattice.CovolumeLe},
\lean{New.TwistedLattice.volume}]\label{setup:def:index}
Let $y = (y_1, \dots, y_n)$ be a family in $V$.
\begin{itemize}
\item $y$ has \emph{index at most $\delta$} if there are $\delta_\ell \ge 0$ with
  $\sum_{\ell \in S}\delta_\ell\log\ell \le \delta$ for all finite sets $S$ of primes, such that
  for every $\iota \mid \ell$ some $n \times n$ matrix $u$ over $\cO_\iota$ with
  $\norm{\iota(\det u)}_\ell \ge \ell^{-\delta_\ell}$ has all combinations $\sum_k u_{ik} y_k$
  integral at $\iota$. The matrix does not depend on the twist.
\item $y$ has \emph{covolume at most $c$} if
  $\abs{\det(\pair{\sigma\Theta_1 y_i, \sigma\Theta_1 y_k})_{i,k}} \le e^{c}$ for every
  $\sigma \in \Aut(\C)$.
\end{itemize}
If $V$ is a finite-dimensional inner product space, put $Z(\xi) = \det(\Theta_\xi^*\Theta_\xi)$,
the Gram determinant of the images under $\Theta_\xi$ of an orthonormal basis of $V$.
\end{definition}

The index measures how far the family is from integral, at the same price for all twists; the
covolume measures its size, for the trivial twist only. By Hadamard's inequality,
if $\norm{\Theta_1 y_i}_\sigma^2 \le e^h$ for all $i$ and $\sigma$, then $y$ has covolume at most
$nh$ (\lean{New.TwistedLattice.covolumeLe_of_normSqAt_le}).

\begin{theorem}[A twist for each prime; \lean{New.TwistedLattice.covolume_add_index_nonneg}]
\label{setup:thm:twist}
Let $V$ be a finite-dimensional inner product space with a twisted lattice all of whose maps
$\Theta_\xi$ are injective. Let $y_1, \dots, y_n, z_1, \dots, z_{n'}$ be a basis of $V$, and for
each $\iota \mid \ell$ let $\xi_\iota$ be a twist allowed at $\ell$. Suppose that
\begin{enumerate}[label=(\roman*)]
\item $y$ has covolume at most $c$ and index at most $\delta$;
\item for each twist $\xi$ there is a number $\gamma(\xi) \ne 0$ such that
  $\gamma(\xi)\Theta_\xi z_i$ is integral at every $\iota \mid \ell$ for which $\xi$ is allowed at
  $\ell$;
\item $\norm{\Theta_1 z_i}_\sigma^2 \le e^{b_{\mathrm{short}}}$ for all $i$ and all
  $\sigma \in \Aut(\C)$;
\item the families $\iota \mapsto \abs{\gamma(\xi_\iota)}^{-2}$ and
  $\iota \mapsto Z(1)/Z(\xi_\iota)$ have height at most $b_{\mathrm{denom}}$ and
  $b_{\mathrm{vol}}$.
\end{enumerate}
Then
\[
  c + 2\delta + n'\,(b_{\mathrm{short}} + b_{\mathrm{denom}}) + b_{\mathrm{vol}} \ge 0 .
\]
\end{theorem}

\begin{proof}
Let $w = (y, z)$ and let $\mathrm{Gram}(\xi)$ be the determinant of the matrix
$(\pair{\Theta_\xi w_i, \Theta_\xi w_k})$. If $B$ is the matrix of $w$ in an orthonormal basis of
$V$, then $\mathrm{Gram}(\xi) = Z(\xi)\,\abs{\det B}^2$. So for every $\iota$
\begin{equation}\label{setup:eq:factor}
  \mathrm{Gram}(1) = \Bigl[ \abs{\gamma(\xi_\iota)}^{2n'}\, \mathrm{Gram}(\xi_\iota) \Bigr]
     \cdot \Bigl[ \abs{\gamma(\xi_\iota)}^{-2} \Bigr]^{n'} \cdot \frac{Z(1)}{Z(\xi_\iota)} .
\end{equation}
The left side does not depend on $\iota$, and it is non-zero because $\Theta_1$ is injective.

\emph{The first factor has height at most $2\delta$}
(\lean{New.TwistedLattice.heightLe_pow_mul_det_gram}). Fix $\iota \mid \ell$, write
$\xi = \xi_\iota$, and let $u$ and $u'$ be the matrices of the definition of the index at $\iota$
and at $\bar\iota$. Enlarge them by the scalar block $\gamma(\xi)$ on the $z_i$. The resulting
matrices carry the vectors $\Theta_\xi w_i$ into $U_\xi(\iota)$ and into $U_\xi(\bar\iota)$. As
$\xi$ is allowed at $\ell$, all inner products between the two transformed families are integral
at $\iota$, hence so is the determinant
$\overline{\det u'}\cdot\abs{\gamma(\xi)}^{2n'}\,\mathrm{Gram}(\xi)\cdot\det u$. Since
$\norm{\iota(\det u)}_\ell$ and $\norm{\bar\iota(\det u')}_\ell$ are at least $\ell^{-\delta_\ell}$, we get
$\norm{\iota(\abs{\gamma(\xi)}^{2n'}\mathrm{Gram}(\xi))}_\ell \le \ell^{2\delta_\ell}$; this is
height at most $2\delta$, with $H = 1$ and $\alpha = 0$.

By (iv) and the additivity of heights, \eqref{setup:eq:factor} shows that the constant family
$\mathrm{Gram}(1)$ has height at most $2\delta + n' b_{\mathrm{denom}} + b_{\mathrm{vol}}$. On the
other hand, as the matrix of the inner product is rational, $\sigma\,\mathrm{Gram}(1)$ is the
determinant of the pairings of the vectors $\sigma'\Theta_1 w_i$ with the $\sigma\Theta_1 w_k$, where
$\sigma'$ is $\sigma$ conjugated by complex conjugation; by the Cauchy--Schwarz inequality for determinants its
square is at most the product of the Gram determinants of these two families. By (i), (iii) and
the inequalities of Fischer and Hadamard, each is at most $e^{c + n' b_{\mathrm{short}}}$, so
$\abs{\sigma\,\mathrm{Gram}(1)} \le e^{c + n' b_{\mathrm{short}}}$ for all $\sigma$
(\lean{New.TwistedLattice.CovolumeLe.sumElim},
\lean{New.TwistedLattice.CovolumeLe.norm_map_det_gram_le}). The product formula,
Lemma~\ref{setup:lem:product}, gives the claim.
\end{proof}

\begin{remark}
This is the determinant method of transcendence theory, by which a non-zero algebraic number
cannot be small at all places (compare Laurent's interpolation determinants, Laurent, 1994, and Bost's slope
inequalities, Bost, 2001), with a twist $\xi_\iota$ for each $\iota$; as
$\mathrm{Gram}(\xi)/Z(\xi)$ does not depend on $\xi$, the price is the height of $Z(1)/Z(\xi_\iota)$.
\end{remark}

\subsection{The lattice of target forms}

The objects named here are defined in Section~\ref{sec:ub-objects}. The \emph{target forms}
$\mathcal{T}_D$ (\lean{New.Data.targetForms}) are a finite-dimensional space of automorphic forms
on the unitary group of a hermitian space $Y_0$ over $E$; put $\dT = \dim \mathcal{T}_D$. Twists are
defined in Definition~\ref{data:def:twist}; in short, a
\emph{twist} $\vartheta$ is an automorphic character of $\U(1)$ of weight $0$ that is tamely
ramified at a finite set $Q$ of places, its \emph{modulus}, disjoint from its conjugate and of norm
prime to $MN$ (\lean{New.Data.TameChar}, \lean{New.Data.Twist}); it is \emph{$m$-fold} if
$\vartheta, \vartheta^2, \dots, \vartheta^m$ all ramify at every place of $Q$
(\lean{New.Data.MFoldTwist}). For an $m$-fold twist, the chain of $m$ theta lifts
$\Theta_\vartheta$ maps $\mathcal{T}_D$ to a space of algebraic modular forms on the definite
unitary group of $Y_m$; such forms are functions on a finite set of classes with values in a
representation with a $\Z$-structure, whence an inner product, integrality at each $\iota$ and
conjugates under $\Aut(\C)$.

\begin{proposition}[\lean{New.Data.targetLattice}]\label{setup:prop:targetlattice}
The maps $\Theta_\vartheta$, with the forms on $Y_m$ that are integral at $\iota$ as integral
vectors, and with $\vartheta$ allowed at $\ell$ if the norm of its modulus is prime to $\ell$, are
a twisted lattice on $\mathcal{T}_D$, with twists the $m$-fold twists. In it,
$\norm{\Theta_1 y}_\sigma^2$ is the squared norm of the conjugate form
(\lean{New.Data.normSqAt_targetLattice}), $Z(\vartheta)$ is the \emph{chain determinant}
$\det(\Theta_\vartheta^*\Theta_\vartheta)$ (\lean{New.Data.volume_targetLattice}), and every
$\Theta_\vartheta$ is injective (\lean{New.Data.injective_lift_targetLattice}).
\end{proposition}

Inner products are integral because the level of the forms on $Y_m$ is so small that all classes
have trivial stabilisers (\lean{New.Data.classAut_lastResidueLevel_eq_bot}); the coordinates for
the trivial twist are the values at class representatives. Injectivity is proved in
Section~\ref{sec:det-chain}.
